% Persian resume - build with: make fa   (XeLaTeX + polyglossia, font: Vazirmatn)
\documentclass[10pt,a4paper]{article}
\usepackage[margin=1.4cm,top=1.1cm,bottom=1.0cm]{geometry}
\usepackage{xcolor}
\usepackage{enumitem}
\usepackage[hidelinks,pdftitle={Hazhir Yousefi - Resume (FA)},pdfauthor={Hazhir Yousefi}]{hyperref}
% bidi (loaded by polyglossia) must come after xcolor/enumitem/hyperref
\usepackage{fontspec}
\usepackage{polyglossia}
\setmainlanguage{farsi}
\setotherlanguage{english}
\setmainfont[Scale=0.9]{Vazirmatn}
\newfontfamily\farsifont[Scale=0.9,Script=Arabic]{Vazirmatn}
\newfontfamily\englishfont[Scale=0.9]{Vazirmatn}
\definecolor{accent}{HTML}{1F4E79}

\pagestyle{empty}
\setlength{\parindent}{0pt}
\setlength{\parskip}{0pt}
\linespread{1.0}
\setlist[itemize]{itemsep=0.5pt,topsep=1.5pt,parsep=0pt,partopsep=0pt}

% Latin run inside Persian text
\newcommand{\en}[1]{\textenglish{#1}}

\renewcommand{\section}[1]{%
  \vspace{6pt}{\bfseries\color{accent}#1}\par\nointerlineskip\vspace{1.5pt}%
  {\color{accent}\hrule height 0.6pt}\vspace{3pt}}

% \entry{title}{organisation}{dates}{location}
\newcommand{\entry}[4]{%
  \vspace{3pt}\textbf{#1}\hfill #3\par
  #2\hfill #4\par\vspace{1pt}}

% \project{name}{tech}{dates}
\newcommand{\project}[3]{%
  \vspace{3pt}\textbf{#1} \textcolor{gray}{|} {\small #2}\hfill #3\par\vspace{1pt}}

\begin{document}

\begin{center}
  {\fontsize{22}{26}\selectfont\bfseries هژیر یوسفی}\\[3pt]
  {\color{accent}\bfseries مهندس نرم‌افزار ارشد \,·\, موبایل، سیستم‌های نهفته و هوش مصنوعی}\\[3pt]
  {\small تهران، ایران \,·\, آمادهٔ همکاری دورکاری و مهاجرت \,·\,
  \en{\href{mailto:Hazhiryousefi2002@gmail.com}{Hazhiryousefi2002@gmail.com}} \,·\,
  \en{+98 993 338 0344}}\\[1pt]
  {\small \en{\href{https://github.com/Hazhir2002}{github.com/Hazhir2002}} \,·\,
  \en{\href{https://linkedin.com/in/hazhir-yousefi-665906232}{linkedin.com/in/hazhir-yousefi-665906232}}}
\end{center}

\section{خلاصه}
مهندس نرم‌افزار ارشد با بیش از ۴ سال سابقه، متخصص در توسعهٔ اپلیکیشن‌های موبایل با \en{Flutter} و سیستم‌های لینوکسی نهفته.
تنها مهندس یک اپلیکیشن پرداخت و یک سیستم‌عامل کیوسک اختصاصی مبتنی بر \en{Raspberry Pi} که روی حدود ۵۰۰ دستگاه فروش خودکار
در سراسر کشور اجرا می‌شود. تجربهٔ توسعهٔ بک‌اند با \en{FastAPI} و \en{Rust} و ساخت محصولات مبتنی بر \en{LLM}
با \en{RAG} و پایپ‌لاین‌های چندعاملی.

\section{سوابق کاری}

\entry{مهندس نرم‌افزار ارشد}{\en{Modaam}}{۲۰۲۴ – اکنون}{تهران (حضوری)}
\begin{itemize}
  \item تنها مهندس اپلیکیشن موبایل \en{Modaam} (\en{iOS} و \en{Android}) برای پرداخت در دستگاه‌های فروش هوشمند، خرید بسته‌های شارژ
        و پیگیری سفارش‌ها؛ مورد استفادهٔ هزاران کاربر در تمام استان‌های ایران.
  \item طراحی و ساخت سیستم‌عامل کیوسک اختصاصی (\en{Raspberry Pi + Flutter}) برای حدود ۵۰۰ دستگاه فروش تنقلات
        (حدود ۱۰٬۰۰۰ تعامل روزانه)، شامل کنترل سخت‌افزار تحویل کالا، پرداخت و پنل تنظیمات فروشنده.
  \item ساخت سیستم به‌روزرسانی از راه دور (\en{OTA}) با \en{Rust} (باینری‌های \en{ARM64} استاتیک با کراس‌کامپایل) برای انتشار
        نسخه‌ها روی کل ناوگان دستگاه‌ها.
  \item سخت‌سازی و بهینه‌سازی لینوکس نهفته (\en{systemd}، قفل‌کردن کیوسک، پشتیبانی از راه دور) برای کار ۲۴/۷؛ عیب‌یابی مشکلات میدانی
        از جمله کرش ناشی از افت ولتاژ، نشت حافظه و \en{race condition} در نمایشگر.
  \item ساخت \en{AXIOM}، دیزاین‌سیستم اختصاصی \en{Flutter} مشترک بین رابط کیوسک و پنل مدیریت (راست‌به‌چپ فارسی)؛
        بسته‌بندی بک‌اند \en{FastAPI} به‌صورت باینری \en{ARM64} با کراس‌کامپایل \en{Docker}.
  \item مسئولیت کامل معماری، تست و استقرار، بدون تیم \en{QA}.
\end{itemize}

\entry{توسعه‌دهندهٔ نرم‌افزار (فرانت‌اند)}{\en{Hotiday}}{فوریه ۲۰۲۴ – مه ۲۰۲۵}{دورکار}
\begin{itemize}
  \item توسعهٔ ماژول‌های وبلاگ، پروفایل، کیف پول و داشبورد (\en{Next.js}، \en{Tailwind CSS}) از روی طراحی \en{Figma} با اتصال به
        \en{REST API} جنگو؛ بهبود \en{SEO} و عملکرد صفحات و تحویل کتابخانهٔ کامپوننت‌های قابل‌استفادهٔ مجدد.
\end{itemize}

\entry{توسعه‌دهندهٔ فول‌استک}{\en{Sazmanyar}}{مارس ۲۰۲۴ – ژوئن ۲۰۲۴}{تهران}
\begin{itemize}
  \item ساخت سیستم مدیریت انبار از صفر (\en{.NET Core}، \en{React}، \en{SQL Server}) با کنترل دسترسی نقش‌محور و ردیابی لحظه‌ای
        موجودی؛ مورد استفادهٔ چند کسب‌وکار در ایران.
\end{itemize}

\entry{توسعه‌دهندهٔ فول‌استک}{\en{Arvin Rayan System}}{ژوئن ۲۰۲۲ – مارس ۲۰۲۴}{تهران}
\begin{itemize}
  \item توسعه و نوسازی ماژول‌های \en{ERP} (\en{.NET Core}، \en{React}، \en{Entity Framework}) و مهاجرت کدبیس قدیمی
        \en{jQuery / .NET Framework} به معماری مدرن و مقیاس‌پذیر.
  \item حضور در جلسات با مشتری و رهبری ارتقای ابزارهای بک‌اند و فرانت‌اند.
\end{itemize}

\section{پروژه‌ها}

\project{\en{Pettaio} – دستیار هوشمند حیوانات خانگی}{\en{Flutter, FastAPI, LangGraph, Chroma, YOLO}}{۲۰۲۴ – اکنون}
\begin{itemize}
  \item توسعهٔ کامل محصول به‌تنهایی (\en{Flutter} با معماری تمیز؛ بک‌اند \en{FastAPI})؛ هم‌بنیان‌گذاران مسئول بخش کسب‌وکار هستند.
  \item پایپ‌لاین چندعاملی (بررسی ایمنی، مسیریاب، متخصص‌های موازی و ترکیب نهایی پاسخ) با چت استریم (\en{SSE}) و حافظهٔ بلندمدت حیوان خانگی.
  \item \en{RAG} با \en{Chroma} و بازرتبه‌بندی \en{cross-encoder}، چارچوب ارزیابی \en{LLM-as-judge} و مدل \en{YOLO} آموزش‌دیده برای تحلیل
        رفتار حیوان؛ \en{CI/CD} روی سرور \en{Hetzner} (\en{Docker Compose}، \en{Nginx}، \en{GitHub Actions}).
\end{itemize}

\project{\en{Tenkai} – اپلیکیشن ژورنال و تناسب اندام}{\en{Flutter, Drift, Rust, Axum, PostgreSQL}}{پروژهٔ شخصی}
\begin{itemize}
  \item اپلیکیشن \en{iOS} با ژورنال صوتی، ردیاب کتاب و ثبت تمرین (\en{hooks\_riverpod}، \en{Drift}).
  \item بک‌اند \en{Rust} (\en{Axum}، \en{Diesel}، \en{PostgreSQL}) با احراز هویت \en{JWT}؛ بیلد و استقرار خودکار با \en{GitHub Actions} روی \en{GHCR}.
\end{itemize}

\section{تحصیلات}
\entry{کارشناسی مهندسی کامپیوتر (نرم‌افزار)}{دانشگاه علم و صنعت ایران}{سپتامبر ۲۰۲۰ – سپتامبر ۲۰۲۵}{تهران}
\begin{itemize}
  \item پایان‌نامه: فرانت‌اند «تاوان»، سامانهٔ برچسب‌گذاری هوشمند اسناد فارسی، برای یک کارفرمای صنعتی واقعی.
\end{itemize}

\section{مهارت‌ها}
\begin{tabular}{@{}p{2.8cm}p{14.4cm}@{}}
  \textbf{موبایل} & \en{Flutter, Dart, BLoC, Riverpod, GoRouter, Drift, Clean Architecture} \\
  \textbf{سیستم‌های نهفته} & \en{Raspberry Pi, embedded Linux, systemd, Rust, OTA updates, Docker cross-compilation} \\
  \textbf{بک‌اند و وب} & \en{Python, FastAPI, Rust (Axum, Diesel), PostgreSQL, SQLite, Django, .NET Core, Next.js, React} \\
  \textbf{هوش مصنوعی} & \en{LLM apps (OpenAI, Gemini), RAG, LangGraph, PydanticAI, Chroma, YOLO, HuggingFace} \\
  \textbf{دواپس و ابزار} & \en{Docker, GitHub Actions, Nginx, GHCR, AWS, Git, Pytest, Figma} \\
  \textbf{زبان‌ها} & فارسی (مادری)، کردی (مادری)، انگلیسی (مسلط، سطح حرفه‌ای)
\end{tabular}

\end{document}
